\section{Discussion}\label{sec:discussion}

\subsection{What the results say}

Three readings of the results bear directly on policy questions the AI-fiscal
conversation is starting to ask. First, the fiscal and household questions
decouple. Capital-side tax design (the realization margin) moves the budget
by hundreds of billions while leaving poverty untouched; the wage-side
distribution moves poverty by \genRapidPovSpanPp pp while net revenue stays
positive throughout. A revenue estimate, however good, is not a household
forecast. Second, the automatic stabilizers are already in the scenario
arithmetic: transfers absorb \genTransferAbsorbPct\% of the gross-tax swing
across the wage-spread variants, which means deficit projections that stop
at the tax code overstate how much the wage-distribution assumption matters
fiscally --- and understate the program-side costs when workers lose. Third,
the state layer is not a rounding error: \genRapidPropState B under
proportional Rapid growth, doubling across wage-spread variants,
concentrated in a handful of states, and available only to states whose
codes reach capital income at all.

\subsection{Limitations}

\textbf{No corporate sector.} PolicyEngine-US models households; the
corporate wedge in Section~\ref{sec:reconciliation} is borrowed from the
Budget Lab, not modelled.

\textbf{Static.} No labor supply, saving, avoidance or enforcement
responses on either side of the comparison.

\textbf{Apportionment is narrower than theirs.} The incremental capital flow
follows existing realized capital income; their SCF-wealth apportionment
reaches households that hold wealth but realize nothing. By their own
account the narrower base is more top-concentrated, so this choice
concentrates the shock more than theirs. The model carries household
balance-sheet variables (\texttt{net\_worth}, \texttt{stock\_assets}), so an
asset-based variant is feasible and is identified future work; the open
modelling question is which person and which income type receives the new
flow in a household that currently realizes none.

\textbf{Pass-through split is coarser than theirs.} Their active/passive-%
aware rule is approximated by a single rule because the microdata carries no
active/passive flag, under-assigning the passive portion to capital.

\textbf{National shares versus modelled shares.} Scenario calibration uses
national factor shares ($\theta^0_L = 0.5550$); observed taxable microdata
flows are more labor-weighted (\$\genModelledLaborShareT T labor against
\$\genModelledCapitalShareT T capital), because survey and tax data miss
unrealized gains, retained earnings and imputed returns. Growth rates from
national accounts are applied to observed taxable flows --- the same choice
the Budget Lab makes with the PUF --- and both baseline shares are reported
so the gap is visible rather than assumed away. One implication is worth
stating for the headline tilt result: under Rapid, the as-forecast scenario
raises modelled market income by \genRapidPropMarketDelta B while the
shares-fixed counterfactual raises it by \genRapidFixedMarketDelta B,
because much of the forecast capital growth accrues to income the taxable
microdata does not observe. The tilt's revenue cost therefore combines a
composition effect with growth landing outside the modelled base; the
realization sweep of Section~\ref{sec:realization} is the sensitivity that
spans this margin.

\textbf{No sampling uncertainty.} Estimates carry survey-sampling and
imputation error that is not quantified here. The build-to-build comparison
of Section~\ref{sec:stability} bounds sensitivity to data calibration ---
material for sub-national and program-level cells --- but not sampling
error; small cells such as individual state takes and single-program
responses should be read to the nearest billion, not the decimal shown.

\textbf{Anchored poverty, unequivalized Gini.} SPM thresholds do not respond
to the simulated distribution, so a fully relative poverty line would rise
with five years of above-baseline growth and show larger increases. The
Gini is household-level and unequivalized: changes are meaningful, levels
are not comparable to published equivalized series. Poverty levels sit above
the published Census SPM rate \citep{census2025} --- a known calibration gap
--- so every poverty figure is a change against a same-year baseline on
identical thresholds.

\textbf{No health-program response.} Medicaid, CHIP and ACA subsidies are
excluded from net-income accounting in these runs by the model's default
switch; the transfer responses reported here would not shrink, and could
grow, with that channel on --- but this paper does not measure it.

\subsection{Next steps}

The wage-inequality parameter $\lambda$ is a stylization with, as the Budget
Lab says plainly, little evidence behind its magnitude. The natural
replacement is measured occupation-level exposure
\citep{felten2021, eloundou2024}: distributing the labor shock by exposure
scores mapped to CPS occupations would make the implied $\lambda$ an output
rather than an input --- the design of the UK companion paper
\citep{ahmadi2026}, and the ``S1'' expansion the Budget Lab's methodology
names as planned future work. Asset-based capital apportionment and a
policy-response layer (recycling the net revenue gain into transfers to test
whether the receipts can hold the bottom of the distribution harmless) are
the other two identified extensions.

\subsection{Conclusion}

Whether AI pays for itself fiscally depends on where its income lands ---
between factors, across the wage distribution, and on or off the individual
tax base. Running forecaster-calibrated scenarios through the full United
States tax and transfer system, the tilt toward capital keeps
\genKeptMatchedRapidPct\% of the revenue that balanced growth would deliver
on the federal basis and \genKeptHouseholdRapidPct\% on the household basis
(states and transfers, no corporate wedge), and it turns growth's poverty
gain into a marginal loss; the wage-distribution assumption nobody can yet
pin down flips the household verdict while the revenue verdict stands; and
the channels a tax-only model omits --- automatic transfers and state codes
--- carry a material share of the response. Where this model and
the Budget Lab's overlap, they agree closely; where they diverge, the
anatomy is identifiable and mostly definitional. All of it is open:
anyone can rerun these scenarios, or score a different future against them.
